% !TeX root = ../main.tex

% 中文摘要
\begin{abstract}
  交通流量预测是智能交通系统中的核心任务之一，能够为交通管理和城市规划提供关键的决策支持。然而，随着城市的不断发展，交通网络日益复杂，规模呈指数级增长，传统的交通流量预测方法面临着许多新的挑战。主要集中在下述方面：(1)交通流量具有强烈的时空依赖性，如何进行时空关联捕获；(2)如何进行多源异质数据融合；(3)如何解决计算复杂度过高的问题等。

为了应对这些挑战，本研究创新性地提出了一种基于时空注意力机制和自适应图卷积网络的联合模型。

时空注意力机制能够自适应地学习时间和空间上的重要特征，帮助模型聚焦于影响交通流量的关键时段和关键区域。其中，时间注意力机制通过为不同时间步的输入赋予不同权重，使模型能够关注交通流量波动较大的时间段，这一机制能显著增强模型捕捉非线性趋势和长期依赖关系的能力。与此同时，空间注意力机制的引入，有效缓解了图卷积的过平滑问题。该机制通过动态调整每个节点与其邻居节点之间的相互影响，空间注意力机制能够避免信息的过度平均化，增强模型对局部特征的敏感性，从而提升预测精度。

其次，在图卷积网络的优化方面，本研究采用了自适应图网络技术，使模型能够学习并反映节点间复杂的潜在关系。这一改进不仅简化了计算流程，而且在处理大规模交通网络数据时显著降低了计算成本，避免了内存溢出和长时间运行的问题，从而提升了模型的实用性。

此外，本研究还在模型中尝试了多源信息的融合，将外部相关数据（如天气条件、POI信息等）纳入考虑。这一举措进一步提升了模型的预测精度，证明了多信息融合在交通流量预测中的有效性。

最后，基于SZ和LA等多源数据集的实验结果表明，本文提出的时空注意力与自适应图卷积联合框架模型在多种复杂场景下的交通流量预测任务中，相较主流基准模型展现出显著优越性。具体而言，在核心性能指标方面，本文提出的模型在LA数据集上预测平均绝对误差（MAE）和均方根误差（RMSE）分别较对比模型下降至少40.70\%和33.08\%。相较于传统图卷积框架，本文采用的自适应图网络学习节点学习机制能有效的缩短模型训练时间。

综上，本文从时空注意力机制、多信息融合、节点自适应学习三个维度提出了一种新的交通流量预测模型。该模型可为新时代智慧交通管理系统提供重要支撑，为优化交通资源配置、缓解交通拥堵、提升城市交通效率提供技术保障。同时，基于自身灵活的适应性和高效的计算效率，本文提出的新模型在处理大规模交通网络数据是具有明显的优势，具有较强的行业应用推广潜力。
\end{abstract}



% 英文摘要
\begin{abstract*}
  Traffic flow prediction is one of the core tasks in intelligent transportation systems, providing key decision support for traffic management and urban planning. However, with the continuous development of cities, the transportation network has become increasingly complex, and its scale has grown exponentially. Traditional traffic flow prediction methods face many new challenges, mainly focusing on the following aspects: 1) the strong spatiotemporal dependencies of traffic flow, how to capture these spatiotemporal correlations; 2) how to fuse multi-source heterogeneous data; 3) how to address the high computational complexity of the problem.

To address these challenges, this study innovatively proposes a joint model based on a spatiotemporal attention mechanism and an adaptive graph convolutional network. The spatiotemporal attention mechanism can adaptively learn important features in both time and space, helping the model focus on key periods and areas that affect traffic flow. The temporal attention mechanism assigns different weights to inputs at different time steps, enabling the model to pay attention to periods of significant traffic flow fluctuations. This mechanism significantly enhances the model's ability to capture nonlinear trends and long-term dependencies. At the same time, the introduction of the spatial attention mechanism effectively alleviates the over-smoothing problem in graph convolution. By dynamically adjusting the interactions between each node and its neighboring nodes, the spatial attention mechanism prevents excessive averaging of information, increases the model's sensitivity to local features, and improves prediction accuracy.

Additionally, in optimizing the graph convolutional network, this study employs an adaptive graph network technique that allows the model to learn and reflect the complex underlying relationships between nodes. This improvement not only simplifies the computational process but also significantly reduces the computational cost when handling large-scale traffic network data, avoiding memory overflow and long processing times, thereby enhancing the model's practicality.

Moreover, this study attempts to incorporate multi-source information into the model by considering external relevant data such as weather conditions and Points of Interest (POI) information. This approach further improves the model's prediction accuracy, demonstrating the effectiveness of multi-source information fusion in traffic flow prediction.

Finally, experimental results based on datasets such as SZ and LA show that the proposed spatiotemporal attention and adaptive graph convolution joint framework model demonstrates significant superiority over mainstream baseline models in traffic flow prediction tasks in various complex scenarios. Specifically, in terms of key performance indicators, The model proposed in this paper achieves at least a 40.70\% and 33.08\% reduction in Mean Absolute Error (MAE) and Root Mean Square Error (RMSE), respectively, compared to the baseline models on the LA dataset. Compared to traditional graph convolution frameworks, the adaptive graph network learning mechanism used in this study can effectively shorten the model's training time.

In conclusion, this study presents a new traffic flow prediction model from the perspectives of temporal attention mechanism, multi-information fusion, and adaptive node learning. This model provides essential support for the new era of smart transportation management systems and offers technical assurance for optimizing traffic resource allocation, alleviating traffic congestion, and improving urban traffic efficiency. Moreover, due to its flexible adaptability and efficient computational performance, the proposed model has significant advantages in handling large-scale traffic network data and demonstrates strong potential for industry application.
\end{abstract*}